% GENERATED FILE. Edit canonical lessons, then run npm run generate:books.
% canonical-source-hash: fnv1a64:ef8184e4fd309678
% canonical-lessons: HI-C06-numbers-1-5, HI-C06-tin-char-history, HI-C06-paanch-nasal

\chapter{Numbers One to Five}
\label{ch:hi-numbers-1-5}

\begin{chapteropening}
\textbf{By the end of this chapter:} \emph{I can count from one to five in Hindi and hear the older Sanskrit shapes still sitting inside \hi{तीन}, \hi{चार} and \hi{पाँच}.}

Say \hi{पाँच} beside Sanskrit pañca and point to the moon-dot that is all that survives of the lost nasal consonant.
\end{chapteropening}

\section[\hi{एक} \hi{दो} \hi{तीन} \hi{चार} \hi{पाँच}]{\hi{एक}, \hi{दो}, \hi{तीन}, \hi{चार}, \hi{पाँच} — the five forms}

\label{lesson:HI-C06-numbers-1-5}



\begin{quote}
Five words, and a story about how languages age.
\end{quote}

\subsection*{You'll want to know — The five}
\noindent
\begin{tabularx}{\linewidth}{@{}>{\raggedright\arraybackslash}X>{\raggedright\arraybackslash}X>{\raggedright\arraybackslash}X@{}}
\toprule
\textbf{} & \textbf{Hindi} & \textbf{said} \\
\midrule
1 & \textbf{\hi{एक}} & \emph{ek} \\
2 & \textbf{\hi{दो}} & \emph{do} \\
3 & \textbf{\hi{तीन}} & \emph{tīn} \\
4 & \textbf{\hi{चार}} & \emph{chār} \\
5 & \textbf{\hi{पाँच}} & \emph{pāṁch} \\
\bottomrule
\end{tabularx}

If you did the writing track you'll recognise some pieces: \textbf{\hi{क}} (Lesson 2), \textbf{\hi{त}} (Lesson 2), \textbf{\hi{न}} (Lesson 1), \textbf{\hi{र}} (Lesson 4), and the \textbf{\hi{ा}} mātrā (Lesson 3).

Plenty here you haven't been taught to draw yet — \textbf{\hi{ए}}, \textbf{\hi{च}}, \textbf{\hi{द}}, \textbf{\hi{प}}, the \textbf{\hi{ी}} and \textbf{\hi{ो}} mātrās, and the \textbf{\hi{ँ}}. Read them now; you'll draw them when the writing track reaches them.

The \textbf{\hi{ँ}} over \emph{pāṁch} is a \emph{chandrabindu}, \textquotedblleft{}moon-dot\textquotedblright{} — it nasalises the vowel, the way French does in \emph{bon}.

\subsection*{Your turn}
\begin{itemize}
\raggedright
  \item \emph{Say it:} \textquotedblleft{}ek, do, tīn, chār, pāṁch\textquotedblright{}
  \item \emph{Point to:} the familiar \hi{क}, \hi{त}, \hi{न}, \hi{र} and \hi{ा} inside the five words
  \item \emph{Trace it:} the \textbf{\hi{ँ}} moon-dot over \emph{pāṁch}
\end{itemize}

\begin{tcolorbox}[breakable,colback=teal!4,colframe=teal!35!black,title={Before you move on}]
Count to five in Hindi. (\emph{Ek, do, tīn, chār, pāṁch}.) Which mark nasalises the vowel in \textbf{\hi{पाँच}}? (The \textbf{chandrabindu}, \textquotedblleft{}moon-dot.\textquotedblright{}) Do you need to draw every unfamiliar letter today? (\textbf{No} — read them now and draw them when their stroke lessons arrive.) Next: watch the Sanskrit forms wear smooth.
\end{tcolorbox}
\section[\hi{तीन}, \hi{चार}]{\hi{तीन} and \hi{चार} — weight moves from consonants to vowels}

\label{lesson:HI-C06-tin-char-history}



\begin{quote}
The five modern forms did not jump straight from Sanskrit. The Prakrits did most of the wearing-down in between.
\end{quote}

\subsection*{You'll want to know — Put the stages side by side}
\noindent
\begin{tabularx}{\linewidth}{@{}>{\raggedright\arraybackslash}X>{\raggedright\arraybackslash}X>{\raggedright\arraybackslash}X@{}}
\toprule
\textbf{Sanskrit} & \textbf{$\to$ Prakrit $\to$} & \textbf{Hindi} \\
\midrule
\emph{éka-} & \emph{ekka} & \textbf{ek} \\
\emph{dvá-} & \emph{do} & \textbf{do} \\
\textbf{neuter \emph{tr\'{\={\i}}ṇi}} & \emph{tiṇṇi} & \textbf{tīn} \\
\textbf{neuter \emph{catv\'{\={a}}ri}} & \emph{cattāri} & \textbf{chār} \\
\emph{páñca} & \emph{pañca} & \textbf{pāṁch} \\
\bottomrule
\end{tabularx}

Three and four come from the \textbf{neuter} forms. Sanskrit also had masculine \emph{tráyaḥ}, but Hindi \emph{tīn} continues \emph{tr\'{\={\i}}ṇi}. The same choice gives \emph{catv\'{\={a}}ri} $\to$ \emph{chār}.

\subsection*{You'll want to know — The two-step weight trade}
1. \emph{\textbf{tr\'{\={\i}}ṇi\emph{ $\to$ }tiṇṇi}}: \emph{r} is lost and \textbf{ṇ doubles} to make up the weight. Long \emph{ī} shortens because the doubled consonant already makes the syllable heavy. 2. \emph{\textbf{tiṇṇi\emph{ $\to$ }tīn}}: the double simplifies and the \textbf{vowel lengthens} to pay for the lost consonant weight.

The vowel ends long again after being short in the middle. \emph{Ekka} $\to$ \emph{ek} shows the second half of the same trade.

This is not decay in a bad sense; it is what heavy use does over two thousand years. The same ancient number wore down on another continent:

\noindent
\begin{tabularx}{\linewidth}{@{}>{\raggedright\arraybackslash}X>{\raggedright\arraybackslash}X>{\raggedright\arraybackslash}X@{}}
\toprule
\textbf{} & \textbf{ancient} & \textbf{modern} \\
\midrule
\textbf{Latin $\to$ Spanish} & \emph{quattuor} & \emph{cuatro} \\
\textbf{Sanskrit $\to$ Hindi} & \emph{catv\'{\={a}}ri} & \emph{chār} \\
\bottomrule
\end{tabularx}

Same PIE word, two continents, two ordinary histories.

\subsection*{Your turn}
\begin{itemize}
\raggedright
  \item \emph{Say it:} \textquotedblleft{}tr\'{\={\i}}ṇi … ti\textbf{ṇṇ}i … t\textbf{ī}n\textquotedblright{}
  \item \emph{Say it:} \textquotedblleft{}catv\'{\={a}}ri … cattāri … \textbf{chār}\textquotedblright{}
  \item \emph{Compare:} \textquotedblleft{}quattuor $\to$ cuatro · catv\'{\={a}}ri $\to$ chār\textquotedblright{}
\end{itemize}

\begin{tcolorbox}[breakable,colback=teal!4,colframe=teal!35!black,title={Before you move on}]
Which Sanskrit forms underlie \emph{tīn} and \emph{chār}? (\textbf{Neuter} \emph{tr\'{\={\i}}ṇi} and \emph{catv\'{\={a}}ri}.) What sits between Sanskrit and Hindi? (The \textbf{Prakrits}.) In \emph{tr\'{\={\i}}ṇi} $\to$ \emph{tiṇṇi} $\to$ \emph{tīn}, what trades places? (The lost \emph{r} is paid for by a \textbf{double consonant}, then the simplified double by a \textbf{long vowel}.) Next: follow the nasal that stayed inside \emph{pāṁch}.
\end{tcolorbox}
\section[pāṁch]{\hi{पाँच} — the nasal moved into the vowel}

\label{lesson:HI-C06-paanch-nasal}



\begin{quote}
\emph{Tīn} moved weight between consonants and vowels. \emph{Pāṁch} moves a whole nasal sound into the vowel itself.
\end{quote}

\subsection*{You'll want to know — The nasal that stayed}
Sanskrit \emph{\textbf{pañca}} $\to$ Hindi \emph{\textbf{pāṁch}}. The \textbf{n} did not vanish; Sanskrit's \emph{ñ} consonant became Hindi's nasalised \emph{ā}. The \textbf{\hi{ँ}} chandrabindu records that nasality above the word.

Say \emph{pāṁch} and \emph{pañca} one after the other. You can hear the nasal migrating backward from its own consonant into the vowel.

\subsection*{Your turn}
\begin{itemize}
\raggedright
  \item \emph{Say it:} \textquotedblleft{}pañca … pāṁch\textquotedblright{}
  \item \emph{Hum:} hold the nasal through the vowel in \emph{pāṁch}
  \item \emph{Point to:} \textbf{\hi{ँ}}, the moon-dot that records it
\end{itemize}

\begin{tcolorbox}[breakable,colback=teal!4,colframe=teal!35!black,title={Before you move on}]
Did the nasal in \emph{pañca} disappear? (\textbf{No} — it moved into the vowel.) What records it in \textbf{\hi{पाँच}}? (The \textbf{chandrabindu}.) Count one to five once more. (\emph{Ek, do, tīn, chār, pāṁch}.)
\end{tcolorbox}
